\par\noindent\SourceLeaf{47}%
\hspace*{\parindent}Zunächst ist klar, dass das mit $\cos 0\varphi = 1$ multiplizierte Glied, d.\,h. das Anfangsglied der neuen Reihe, allein aus der Auflösung des Anfangsgliedes in eine Summe hervorgeht. Bei den übrigen Gliedern enthalten beide Teile der jeweils entstehenden Summe den Kosinus eines geraden Vielfachen von $\varphi$. Das Anfangsglied der neuen Reihe ist also
\[
\frac{1}{2^{2m-2}} . \frac{1}{2}\,{}^{2m-1}\mathfrak{B}^{(m)}
= \frac{1}{2^{2m-1}}\,{}^{2m-1}\mathfrak{B}^{(m)}.
\]

\par\noindent Dieser Binomialkoeffizient ist der spätere der mittleren Koeffizienten seiner Potenz. Addiert man zu ihm den anderen mittleren Koeffizienten, dem er gleich ist, so erhält man den Koeffizienten mit demselben Index für die nächsthöhere Potenz. Es gilt daher
\[
{}^{2m-1}\mathfrak{B}^{(m)} .  2 = {}^{2m}\mathfrak{B}^{(m)},
\]
und der Ausdruck lässt sich so darstellen. Das Anfangsglied ist
\[
= \frac{1}{2^{2m-1}}\frac{{}^{2m}\mathfrak{B}^{(m)}}{2}.
\]

\par\noindent Zur Bestimmung des allgemeinen Gesetzes für die übrigen Glieder nehme man aus\hfill\mbox{}\linebreak $\cos\varphi^{2m-1} . \cos\varphi$
das allgemeine Glied $r^{\text{tus}}$
\[
\frac{1}{2^{2m-2}}\,{}^{2m-1}\mathfrak{B}^{(m+r)}\cos(2r+1)\varphi . \cos\varphi.
\]
Da
\[
\cos(2r+1)\varphi . \cos\varphi = \tfrac{1}{2}\cos 2r\varphi + \tfrac{1}{2}\cos(2r+2)\varphi,
\]
lässt sich dieses Glied in die folgenden Teile auflösen
\[
\frac{1}{2^{2m-2}} . \frac{1}{2} . {}^{2m-1}\mathfrak{B}^{(m+r)}\cos 2r\varphi
+ \frac{1}{2^{2m-2}} . \frac{1}{2} . {}^{2m-1}\mathfrak{B}^{(m+r)}\cos(2r+2)\varphi.
\]

Daraus ersieht man, dass bei dieser Umformung aus jedem ungeraden Vielfachen zwei gerade hervorgehen, das nächsthöhere und das nächstniedrigere. So treten in der neuen Reihe nur Kosinus gerader Vielfacher auf. Jeder erhält einen Koeffizienten, der sich aus zwei benachbarten Koeffizienten der angenommenen Reihe zusammensetzt. Aus der Umformung des Gliedes $r^{\text{to}}$ von
\[
\cos\varphi^{2m-1} . \cos\varphi
\]
geht nämlich hervor, dass sowohl dessen bereits gefundener erster Teil
\[
\frac{1}{2^{2m-1}}\,{}^{2m-1}\mathfrak{B}^{(m+r)}\cos 2r\varphi,
\]
